\section*{Chapter \ref{ch:rc:the-valuation-adjustment-desk} --- The Valuation-Adjustment Desk}

\begin{solution}{exo:rc:the-valuation-adjustment-desk:1}
Netting: the new trade's values offset the existing set's on many paths, so the netted positive
exposure falls. Standalone, the trade has its own positive exposure; in the set, it removes more
exposure than it adds.
\end{solution}

\begin{solution}{exo:rc:the-valuation-adjustment-desk:2}
$1\,000\,000/(100\,000\,000\times13.67) = 7.3$ basis points a year.
\end{solution}

\begin{solution}{exo:rc:the-valuation-adjustment-desk:3}
A measure of credit quality (such as a rating), industry, and region.
\end{solution}

\begin{solution}{exo:rc:the-valuation-adjustment-desk:4}
$\mathrm{CVA} = (1-R)\sum_k\E[D\max(\sum_iV_i,0)]\Delta\mathrm{PD}_k$ and
$\max(\sum_iV_i,0) = \sum_iV_i\mathbf 1_{\sum_jV_j>0}$ path by path; take expectations and sum over $i$:
the contributions $a_i$ add up. Equivalently, CVA is homogeneous of degree one in the sizes and
Euler's theorem applies.
\end{solution}

\begin{solution}{exo:rc:the-valuation-adjustment-desk:5}
The incremental charge of a trade depends on what is already in the set, so two offsetting
trades are charged differently depending on which came first. Euler contributions are computed on
the whole set at once and are symmetric in the trades, so they do not depend on history.
\end{solution}

\begin{solution}{exo:rc:the-valuation-adjustment-desk:6}
Basis between the client's spread and the index, jump-to-default of the client (index protection
pays little when one name defaults), \omterm{def:rc:counterparty-exposure:wwr}{wrong-way risk}, and the changes of the exposure itself.
\end{solution}

\begin{solution}{exo:rc:the-valuation-adjustment-desk:7}
USD~2\,385 per basis point; divided by the peer's ten-year \omterm{def:rc:reduced-form-credit:risky}{risky annuity} and a basis point,
about USD~3.55 million of protection.
\end{solution}

\begin{solution}{exo:rc:the-valuation-adjustment-desk:8}
The charge is the adjustment's value at inception; afterwards the desk's P\&L moves with the
client's spread, the exposure's drivers, funding and capital costs, and any unhedged gap. Only a
perfect hedge of every component would keep it flat, and funding and capital cannot be hedged.
\end{solution}

\begin{solution}{pb:rc:the-valuation-adjustment-desk:1}
\textbf{1.} 4.11\%; an annuity of 13.67.
\textbf{2.} 337.5 basis points, from the peer regression with the client mapped to BB,
industrials, US.
\textbf{3.} CVA USD~1\,649\,494 (12.1 basis points), FCA 343\,872 (2.5), KVA 3\,141\,826 (23.0).
\textbf{4.} 37.6 basis points a year.
\textbf{5.} KVA: an unrated client gets the 7\% CVA risk weight and a twenty-year swap ties up
capital for decades.
\textbf{6.} The CVA, 32\% of the total; funding and capital cannot be hedged in the market.
\textbf{7.} USD~2\,385 per basis point; about USD~3.55 million of ten-year peer protection.
\textbf{8.} Basis to the client's own credit, its jump-to-default, the exposure's growth beyond
ten years, and \omterm{def:rc:rates-risk:crossgamma}{cross-gamma}.
\textbf{9.} Rarely for an unrated corporate; if available it would cost about the CVA plus the
seller's margin and capital.
\textbf{10.} With swaps and swaptions sized on the CVA's rate delta and vega, on common
scenarios.
\textbf{11.} A better internal rating, a lower funding spread, a lower \omterm{def:rc:funding-margin-and-capital-adjustments:kva}{hurdle rate}, a
different CVA capital approach, or pricing only the incremental effect on a large existing
relationship.
\textbf{12.} With BBB and the 3\% risk weight the add-on falls to 26.2 basis points (CVA 7.2,
FCA 3.1, KVA 15.8).
\textbf{13.} A CSA, a break clause at ten years, a shorter swap rolled over, or collateral.
\textbf{14.} It may be mispricing the client's credit or not charging capital, and it will
collect the clients whose risk it underprices.
\textbf{15.} Only if the trade would have paid for its risks; walking away from an
underpriced trade avoids a loss.
\textbf{16.} The bank sets the hurdle; the \omterm{def:rc:the-valuation-adjustment-desk:desk}{XVA desk} should charge it to the originating desk so
the trade's economics are visible where the decision is made.
\textbf{17.} Sales credit should be net of the \omterm{def:rc:the-valuation-adjustment-desk:charge}{XVA charge}, so that the salesperson is not paid for
revenue that the bank's risks consume.
\textbf{18.} That the rate includes the cost of its credit, of funding an unsecured position and
of the capital the bank holds, and that a CSA or break clause would lower it.
\textbf{19.} Named result: \emph{the unrated corporate}: an all-in add-on of 37.6 basis points a
year, of which the desk can hedge the CVA's 12.1 basis points (32\%) in the market.
\textbf{20.} The management of counterparty, funding and capital costs that no single desk can
see or hedge alone.
\end{solution}

\begin{solution}{iq:rc:the-valuation-adjustment-desk:1}
Adjustments are properties of \omterm{def:rc:counterparty-exposure:netting}{netting sets} spanning desks; only a central view can price
netting and collateral, avoid double counting, hedge credit centrally, and manage funding and
capital consistently.
\iqlookfor{\omterm{def:rc:counterparty-exposure:netting}{netting sets} across desks and central hedging.}
\end{solution}

\begin{solution}{iq:rc:the-valuation-adjustment-desk:2}
Each trade's contribution is the derivative of the total with respect to scaling that trade;
for a function homogeneous of degree one, contributions add up to the total (Euler's theorem).
CVA with fixed credit is homogeneous; capital formulas with floors or caps may not be.
\iqlookfor{the derivative definition and the homogeneity condition.}
\end{solution}

\begin{solution}{iq:rc:the-valuation-adjustment-desk:3}
Build a proxy curve from liquid peers by rating, sector and region (regression or buckets), or
map to a related name; validate against bonds or loans if any; hedge with index or peer
protection and hold reserves for the basis.
\iqlookfor{proxy construction and residual basis.}
\end{solution}

\begin{solution}{iq:rc:the-valuation-adjustment-desk:4}
Keep each \omterm{def:rc:counterparty-exposure:netting}{netting set}'s path values on the common scenarios in memory; price the new trade on
the same paths (closed forms or precomputed grids); recompute the netted profiles and
adjustments incrementally; parallelise; fall back to approximations for exotic trades.
\iqlookfor{stored paths and incremental recomputation.}
\end{solution}

\begin{solution}{iq:rc:the-valuation-adjustment-desk:5}
Changes in counterparty spreads and exposures net of hedges, funding costs, capital, new-trade
charges and hedge slippage. In a crisis spreads and exposures jump together (\omterm{def:rc:counterparty-exposure:wwr}{wrong-way risk}),
proxies decouple from clients, hedges become expensive, and CVA capital hedges can themselves
be charged.
\iqlookfor{the P\&L drivers and crisis dynamics.}
\end{solution}

\begin{solution}{iq:rc:the-valuation-adjustment-desk:6}
In full, the bank loses price-sensitive clients to banks with cheaper assumptions and keeps
those whose risk others price higher (adverse selection); in part, it subsidises trades. Banks
pass most of it, compete on assumptions and relationship value, and must be consistent.
\iqlookfor{competition, adverse selection and consistency.}
\end{solution}
